\section{Experimental Setup}

The parameters of the empirical investigation are detailed here, encompassing dataset characteristics, partitioning procedures, training hyperparameters, execution environment, and performance indices.

\subsection{Datasets}

Evaluation is performed using three standard public video datasets. The UMN dataset represents the primary source for training and testing routines. In contrast, the UCSD Pedestrian and CUHK Avenue datasets are used solely for zero-shot evaluation of domain generalisation. A comparison of these sources is outlined in Table~\ref{tab:datasets}.

\begin{table*}[!t]
\centering
\caption{Comparison of the three crowd datasets used in this study. The UMN source is used for training and test-set verification; the UCSD and CUHK sets are reserved for zero-shot cross-dataset evaluation.}
\label{tab:datasets}
\footnotesize
\begin{tabular}{llllll}
\toprule
\textbf{Dataset} & \textbf{Scene Type} & \textbf{Videos} & \textbf{Approx.\ Frames} & \textbf{Resolution} & \textbf{Usage} \\
\midrule
UMN              & Indoor/Outdoor Crowd & 11 sequences  & ${\sim}6{,}165$  & $320 \times 240$ & Training \& Testing \\
UCSD Ped2        & Pedestrian Walkway   & 12 train + 12 test & ${\sim}4{,}560$  & $360 \times 240$ & Cross-Dataset Eval. \\
CUHK Avenue      & Campus Entrance      & 16 train + 21 test & ${\sim}30{,}652$ & $640 \times 360$ & Cross-Dataset Eval. \\
\bottomrule
\end{tabular}
\end{table*}

\subsubsection{UMN Crowd Activity Dataset}
The UMN dataset \cite{mehran2009abnormal} is comprised of 11 video clips captured across three locations: an indoor lobby (Scene 1, 3 clips), an outdoor lawn (Scene 2, 4 clips), and an outdoor plaza (Scene 3, 4 clips). All video recordings were captured with 30~FPS frame rate and $320 \times 240$ resolution. In each video, standard pedestrian flow turns into an episode of panic, where people start running away from each other. The total length of all video sequences is approximately 6{,}165 frames.

Since the exact frame indices of anomalies are not specified, we divided each sequence into training and testing sets by taking the last 30 percent of frames as anomalous and the preceding 70 percent as normal. In reality, such a split would correspond to a case when the panic begins near the end of the recording. The implications of potential label noise from this approximation are analyzed in Section~6.

\subsubsection{UCSD Pedestrian Dataset (Ped2 Subset)}
The UCSD dataset \cite{mahadevan2010anomaly} contains video sequences captured from an elevated, static camera view of a pedestrian sidewalk. The Ped2 subset, consisting of 12 training and 12 test sequences at $360 \times 240$ resolution, is used. Anomaly events in the test sequences are defined by the appearance of non-pedestrian objects (e.g., bicycles, skateboards, and vehicles) that disrupt standard walking patterns. Frame-level annotations are included. To test generalisation capabilities across different views, resolutions, and anomaly categories, the model trained on UMN is evaluated directly on the Ped2 test clips without parameter updates or fine-tuning.

\subsubsection{CUHK Avenue Dataset}
The CUHK Avenue dataset \cite{lu2013abnormal} includes 16 training and 21 testing sequences recorded at a campus entryway with a resolution of $640 \times 360$ pixels. The test sequences contain 47 anomaly events, including running, loitering, throwing objects, and counter-flow walking. In line with the UCSD testing protocol, the CUHK Avenue dataset is used only for zero-shot evaluation to verify the transferability of the nine-feature descriptor and UMN-trained model to different environments.

\subsection{Data Partitioning}

Training and parameter validation are performed using exclusively UMN data. The 11 video sequences are processed using the feature extraction rules (Section~3.4) and the temporal window constructor (Section~3.5) to produce roughly 1{,}100 sequences. To preserve the proportion of normal versus anomalous samples, the dataset is split into three distinct, non-overlapping subsets using stratified random sampling:

\begin{itemize}
  \item \textbf{Training set (70\%):} Used for weights adjustment and model parameter optimization.
  \item \textbf{Validation set (15\%):} Utilized for performance monitoring, learning rate scheduling, and early stopping trigger control.
  \item \textbf{Test set (15\%):} Set aside exclusively for reporting final performance metrics, remaining untouched during model selection and training.
\end{itemize}

Matching the class ratios across the three divisions ensures that metric evaluation is not biased by unbalanced distributions.

\subsection{Training Configuration}

Table~\ref{tab:hyperparams} outlines the hyperparameters and configuration parameters used for optimization.

\begin{table}[H]
\centering
\caption{Hyperparameter configuration for classifier training.}
\label{tab:hyperparams}
\footnotesize
\setlength{\tabcolsep}{4pt}
\begin{tabular}{p{4.2cm}p{3.2cm}}
\toprule
\textbf{Hyperparameter} & \textbf{Value} \\
\midrule
BiLSTM hidden units (per direction) & 128 \\
BiLSTM layers & 2 \\
Attention dimension ($d_a$) & 128 \\
Dropout rate & 0.3 \\
Batch size & 32 \\
Optimiser & Adam ($\beta_1 = 0.9$, $\beta_2 = 0.999$) \\
Initial learning rate & $1 \times 10^{-3}$ \\
LR scheduler & ReduceLROnPlateau (factor $= 0.5$, patience $= 5$) \\
Maximum epochs & 100 \\
Early stopping patience & 15 epochs \\
Loss function & Binary Cross-Entropy with Logits \\
Classification threshold & 0.5 \\
Temporal window length ($L$) & 15 frames \\
Sequence length ($S$) & 10 windows \\
Motion threshold ($\tau$) & 0.5 pixels/frame \\
\bottomrule
\end{tabular}
\end{table}

For the training phase, model parameters are updated using the Adam optimizer. A ReduceLROnPlateau scheduler dynamically manages the learning rate: if the validation loss shows no improvement for 5 consecutive epochs, the learning rate is halved. This helps the network escape local minima. To prevent overfitting, training terminates early if the validation loss does not improve for 15 epochs.

\subsection{Hardware and Software Environment}

Workstation specifications and software packages are listed in Table~\ref{tab:hardware}. A single laptop workstation was utilized for all computational runs. For reproducibility, a fixed seed value of 42 was initialized across all libraries, including Python, NumPy, PyTorch, and CUDA.

\begin{table}[H]
\centering
\caption{Hardware and software execution environment details.}
\label{tab:hardware}
\footnotesize
\setlength{\tabcolsep}{4pt}
\begin{tabular}{p{2.8cm}p{4.9cm}}
\toprule
\textbf{Component} & \textbf{Specification} \\
\midrule
GPU & NVIDIA GeForce RTX 4050 (6 GB VRAM) \\
CPU & Intel Core i7-13650HX \\
RAM & 16 GB DDR5 \\
Operating System & Windows 11 \\
Python & 3.12 \\
PyTorch & 2.0 \\
OpenCV & 4.8 \\
NumPy & 1.26 \\
scikit-learn & 1.3 \\
Random seed & 42 \\
\bottomrule
\end{tabular}
\end{table}

\subsection{Evaluation Metrics}

Five classification metrics are used to measure test-set performance:

\begin{itemize}
  \item \textbf{Accuracy} measures the overall fraction of sequences classified correctly:
  \begin{equation}
    \text{Accuracy} = \frac{\text{TP} + \text{TN}}{\text{TP} + \text{TN} + \text{FP} + \text{FN}}
  \end{equation}

  \item \textbf{Precision} quantifies the proportion of positive predictions that are correct, reflecting the false alarm rate:
  \begin{equation}
    \text{Precision} = \frac{\text{TP}}{\text{TP} + \text{FP}}
  \end{equation}

  \item \textbf{Recall} (sensitivity) measures the proportion of actual stampede events that are successfully detected:
  \begin{equation}
    \text{Recall} = \frac{\text{TP}}{\text{TP} + \text{FN}}
  \end{equation}

  \item \textbf{F1-Score} is the harmonic mean of precision and recall, providing a single balanced metric:
  \begin{equation}
    F_1 = 2 \cdot \frac{\text{Precision} \cdot \text{Recall}}{\text{Precision} + \text{Recall}}
  \end{equation}

  \item \textbf{ROC-AUC} is the area under the Receiver Operating Characteristic curve, which plots the true positive rate against the false positive rate across all classification thresholds. ROC-AUC measures the model's discriminative ability independently of the chosen threshold.
\end{itemize}

To compute precision, recall, and F1-score, binary decisions are generated from the sigmoid output using a threshold of 0.5. For crowd safety installations, high precision is essential: frequent false alarms can lead to operator fatigue, causing genuine alerts to be overlooked.
